\section*{第 \ref{ch:b1:earthquake-source} 章——地震震源}

\begin{solution}{exo:b1:earthquake-source:1}
(a) $\log_{10}(\num{4e17}) = 17.60$，故 $M_w = \tfrac23(17.60 - 9.1) = 5.67$。
(b) $\log_{10}(\num{5.59e22}) = 22.75$，且 $M_w = \tfrac23(22.75 - 9.1) =
9.10$。(c) $M_0 = 10^{1.5\times7.0 + 9.1} = 10^{19.6} = \qty{3.98e19}{N.m}$。
\end{solution}

\begin{solution}{exo:b1:earthquake-source:2}
(a) $T = D/\dot s = \qty{3}{m}/\qty{0.025}{m.\text{年}^{-1}} = \qty{120}{\text{年}}$。
(b) 取 $c = 0.5$，则 $T = D/(c\dot s) = \qty{240}{\text{年}}$：蠕滑承担一半的运动，
而同等滑移量的地震出现频率减半。
\end{solution}

\begin{solution}{exo:b1:earthquake-source:3}
由\cref{prop:b1:earthquake-source:moment-steps}可知，两者矩之比为
$10^{1.5\times(7.4 - 6.2)} = 10^{1.8} = 63.1$：释放与一次震级 7.4 地震相同的矩，约需 63 次
震级 6.2 的地震。
\end{solution}

\begin{solution}{exo:b1:earthquake-source:4}
取 $b = 1$，震级每增加 1，次数减为十分之一：每年震级 5 及以上的地震约 10 次，
震级 6 及以上的约 1 次，震级 7 及以上的约 0.1 次，
即平均每 10 年一次。
\end{solution}

\begin{solution}{exo:b1:earthquake-source:5}
(a) $A = \num{4.3e5} \times \qty{1.5e4}{m} = \qty{6.45e9}{m^2}$，且
$M_0 = \num{3e10}\times\num{6.45e9}\times4 = \qty{7.74e20}{N.m}$；
$\log_{10}M_0 = 20.89$，$M_w = \tfrac23(20.89 - 9.1) = 7.86$。(b) 深度取
\qty{10}{km} 时，面积与矩均除以 1.5，震级因此降低
$\tfrac23\log_{10}1.5 = 0.12$，为 7.74。(c) 第一个结果 7.86 四舍五入为
目录中的 7.9；第二个结果略低。震级对深度并不十分敏感，
面积变化 1.5 倍仅使震级改变 0.12。
\end{solution}

\begin{solution}{exo:b1:earthquake-source:6}
(a) $D = M_0/(\mu A) = \num{5.59e22}/(\num{3e10}\times\num{4e5}\times\num{1.5e5})
= \qty{31.1}{m}$。(b) $\qty{31.1}{m}/\qty{0.083}{m.\text{年}^{-1}} = \qty{374}{\text{年}}$。
(c) 取闭锁比例 0.8，$31.1/(0.8\times0.083) = \qty{468}{\text{年}}$；取
0.5，则为 \qty{749}{\text{年}}。闭锁程度越低，
一次大地震的滑移积累所需时间越长；此处用一个
\qtyproduct{400 x 150}{km} 的矩形、$\mu = \qty{30}{GPa}$ 所得的平均滑移量本身
只是估计，且实际滑移远非均匀。
\end{solution}

\begin{solution}{exo:b1:earthquake-source:7}
(a) \omterm{def:b1:geological-maps:attitude}{走向} $\phi = 329$，\omterm{def:b1:geological-maps:attitude}{倾向} $a = 059$，\omterm{def:b1:geological-maps:attitude}{倾角} $\delta = \ang{8}$。
于是 $\mathbf s = (\sin329,\allowbreak \cos329,\allowbreak 0) =\allowbreak (-0.5150,\allowbreak 0.8572,\allowbreak 0)$，
$\mathbf u = (-\cos8\sin59,\allowbreak -\cos8\cos59,\allowbreak \sin8) =\allowbreak (-0.8488,\allowbreak -0.5100,\allowbreak 0.1392)$，
$\mathbf n = (\sin8\sin59,\allowbreak \sin8\cos59,\allowbreak \cos8) =\allowbreak (0.1193,\allowbreak 0.0717,\allowbreak 0.9903)$，
取 $r = \ang{110}$，
\[
  \mathbf d = \cos110\,\mathbf s + \sin110\,\mathbf u = (-0.6215, -0.7724, 0.1308) .
\]
(b) \omterm{def:b1:earthquake-source:focal-mechanism}{辅助面}的法向为 $\mathbf d$，滑动方向为 $\mathbf n$。
取法向向上，其\omterm{def:b1:geological-maps:attitude}{倾角}为 $\arccos 0.1308 = \ang{82.5}$，\omterm{def:b1:geological-maps:attitude}{倾向}为
$\arctan(-0.6215/-0.7724) = 218.8$，\omterm{def:b1:geological-maps:attitude}{走向}即为 $128.8$；将滑动
$\mathbf n$ 投影到其\omterm{def:b1:geological-maps:attitude}{走向}与上倾矢量上，得\omterm{def:b1:faults-folds-fabrics:fault-elements}{侧伏角}
\ang{87.2}。取整为 $129/83/87$，即已发表解答中的第二个\omterm{def:b1:earthquake-source:focal-mechanism}{节面}。
(c) $\mathbf P = (\mathbf n - \mathbf d)/\sqrt2$ 的方位为 $221$，\omterm{def:b1:field-lab-methods-1:trend-plunge}{倾伏角}
\ang{37}；$\mathbf T = (\mathbf n + \mathbf d)/\sqrt2$ 的方位为 $036$，\omterm{def:b1:field-lab-methods-1:trend-plunge}{倾伏角}
\ang{52}（已发表结果为 $222/38$ 与 $36/52$）。(d) 浅层\omterm{def:b1:earthquake-source:focal-mechanism}{节面}，
\omterm{def:b1:geological-maps:attitude}{倾角} \ang{8}、\omterm{def:b1:geological-maps:attitude}{倾向} $059$，沿\omterm{def:b1:active-margins:slab}{俯冲板片}延伸：为\omterm{def:b1:faults-folds-fabrics:fault-types}{逆断层}，\omterm{def:b1:faults-folds-fabrics:fault-elements}{侧伏角}
为 \ang{110}；又因\omterm{def:b1:geological-maps:attitude}{倾角}小于 \ang{45}，属逆冲，即
俯冲带的板块界面。
\end{solution}

\begin{solution}{exo:b1:earthquake-source:8}
(a) 该面\omterm{def:b1:geological-maps:attitude}{倾向} $135$，故 $\mathbf n = (\sin135, \cos135, 0) =
(0.707, -0.707, 0)$；\omterm{def:b1:faults-folds-fabrics:fault-types}{右行}滑动取 $r = \ang{180}$，$\mathbf d = -\mathbf s
= (-0.707, -0.707, 0)$。朝北，$\mathbf x = (0, 1, 0)$：
$(\mathbf n\cdot\mathbf x)(\mathbf d\cdot\mathbf x) = (-0.707)(-0.707) > 0$，为
挤压。朝东：$(0.707)(-0.707) < 0$，为拉张。朝南为挤压；朝西为拉张。
(b) $\mathbf P = (\mathbf n - \mathbf d)/\sqrt2 = (1, 0, 0)$：\omterm{def:b1:earthquake-source:focal-mechanism}{P 轴}水平，
方位为 $090$--$270$；$\mathbf T = (0, -1, 0)$ 的方位为
$000$--$180$。(c) \omterm{def:b1:geological-maps:attitude}{走向}为 $135$--$315$ 的竖直面，其法向为
$\mathbf d$。配合滑动 $\mathbf n$，位于其西南侧的块体（$\mathbf d$ 指向该侧）
向 $135$ 方向运动；从东北侧块体望向西南，
这一运动位于左侧：为\omterm{def:b1:faults-folds-fabrics:fault-types}{左行}滑动。由
\cref{prop:b1:earthquake-source:auxiliary}可知，此滑动在各方向给出相同的
\omterm{def:b1:earthquake-source:focal-mechanism}{初动}。
\end{solution}

\begin{solution}{exo:b1:earthquake-source:9}
(a) $V = \qty{40}{mm.\text{年}^{-1}}$，在 \qty{10}{km} 处 $v = V/4$，故由
\cref{prop:b1:earthquake-source:profile}(b) 得 $D = \qty{10}{km}$；直接计算，
$D = x/\tan(\pi v/V) = 10/\tan(\pi/4) = \qty{10}{km}$。(b) $v(30) =
(40/\pi)\arctan 3 = \qty{15.9}{mm.\text{年}^{-1}}$。(c) $V/(\pi D) =
\num{0.040}/(\pi\times\num{10000}) = \num{1.27e-6}$ 每年。
\end{solution}

\begin{solution}{exo:b1:earthquake-source:10}
(a) 各区间的矩仍为 $C q^k$，其中 $q = 10^{1.5 - b}$，但此时
$q < 1$：区间的矩随震级增大而减小，且
最底层区间 $k = 0$ 所占份额为 $C/S = (1 - q)/(1 - q^{K+1}) > 1 - q$，恰与
\cref{prop:b1:earthquake-source:largest}成镜像。最小的
地震将占主导，且随着纳入的地震越来越小，
总量会无限增长。(b) $1 - 10^{-0.7} = 0.80$：至少
\qty{80}{\percent}。(c) $\log_{10}N = 8.20 - 1.01\times9 = -0.89$，故 $N = 0.13$
每年，五十年约 6.4 次，为所观测两次的三倍。直线
拟合于震级 6 与 7.5 之间，高估了最大地震；
要么定律在顶端向下弯曲，如断层有限尺寸所暗示的，要么
五十年的样本不足以捕捉数十年才出现一次的事件。两种解释
仅凭两次事件，二者皆有可能。
\end{solution}

\begin{solution}{exo:b1:earthquake-source:11}
(a) 150 年内，位于 $x$ 处的测桩相对断层上测桩的位移为 $150\,v(x) = 150\times(34/\pi)
\arctan(x/15)$ 毫米：在 \qty{1}{km} 处为 \qty{0.11}{m}，
在 \qty{15}{km} 处为 \qty{1.28}{m}，在 \qty{50}{km} 处为 \qty{2.08}{m}，
对\omterm{def:b1:faults-folds-fabrics:fault-types}{右行}断层，自西侧望去向南。(b) 滑移亏损为
$s = \qty{0.034}{m.\text{年}^{-1}}\times\qty{150}{\text{年}} = \qty{5.1}{m}$。地震期间，
测桩的移动为 $(s/\pi)\arctan(D/x)$：在
\qty{1}{km} 处为 \qty{2.44}{m}，在 \qty{15}{km} 处为 \qty{1.28}{m}，在 \qty{50}{km} 处为 \qty{0.47}{m}。
(c) 求和得 $0.11 + 2.44$、$1.28 + 1.28$ 与 $2.08 + 0.47$，均为
\qty{2.55}{m} $= s/2$：地震之后，篱笆在
东侧重新变直，位移为 $s/2$；由对称性，西侧位移为 $-s/2$。
\end{solution}

\begin{solution}{exo:b1:earthquake-source:12}
(a) $\mathbf P = (0, 1, 0)$ 与 $\mathbf T = (1, 0, 0)$。取
$\mathbf n = (\mathbf T + \mathbf P)/\sqrt2 = (0.707, 0.707, 0)$,
$\mathbf d = (\mathbf T - \mathbf P)/\sqrt2 = (0.707, -0.707, 0)$ 满足
$\mathbf P = (\mathbf n - \mathbf d)/\sqrt2$ 与 $\mathbf T = (\mathbf n +
\mathbf d)/\sqrt2$。两矢量均为水平向，故两个\omterm{def:b1:earthquake-source:focal-mechanism}{节面}均
为竖直面：法向为 $\mathbf n$ 的面（朝向 $045$）\omterm{def:b1:geological-maps:attitude}{走向}为 $135$--$315$，
法向为 $\mathbf d$ 的面（朝向 $135$）\omterm{def:b1:geological-maps:attitude}{走向}为 $045$--$225$。(b) 在
\omterm{def:b1:geological-maps:attitude}{走向}为 $315$ 的面上，东北侧的块体沿 $\mathbf d$ 运动，
朝向 $135$：从另一块体看去，向右运动，为\omterm{def:b1:faults-folds-fabrics:fault-types}{右行}走滑。在
\omterm{def:b1:geological-maps:attitude}{走向}为 $045$ 的面上，东南侧的块体沿 $\mathbf n$ 运动，
朝向 $045$：\omterm{def:b1:faults-folds-fabrics:fault-types}{左行}走滑。(c) 新断层的情形为 $\ang{45} - \phi/2 =
\ang{30}$，即新断层与 $\sigma_1$ 成此夹角。若 $\sigma_1$ \omterm{def:b1:geological-maps:attitude}{走向}为 $000$，则断层
\omterm{def:b1:geological-maps:attitude}{走向}分别为 $030$（\omterm{def:b1:faults-folds-fabrics:fault-types}{左行}）与 $330$（\omterm{def:b1:faults-folds-fabrics:fault-types}{右行}），与\omterm{def:b1:earthquake-source:focal-mechanism}{节面}相差 \ang{15}。
\omterm{def:b1:earthquake-source:focal-mechanism}{P 轴}由几何关系固定在距\omterm{def:b1:earthquake-source:focal-mechanism}{节面} \ang{45} 处，无论岩石的摩擦如何，
因此只有断层在与 $\sigma_1$ 成
\ang{45} 处滑动时，\omterm{def:b1:earthquake-source:focal-mechanism}{P 轴}才与之重合。(d) \omterm{def:b1:geological-maps:attitude}{走向}为 $040$ 的断层与\omterm{def:b1:earthquake-source:focal-mechanism}{节面}相差 \ang{5}，
\omterm{def:b1:geological-maps:attitude}{走向} $045$：该\omterm{def:b1:earthquake-source:focal-mechanism}{节面}很可能就是断层，其上的滑动为
\omterm{def:b1:faults-folds-fabrics:fault-types}{左行}。
\end{solution}

\begin{solution}{pb:b1:earthquake-source:1}
\textbf{1.} 间隔分别为 $1290 - 1050 = 240$、$1480 - 1290 = 190$ 与
$1712 - 1480 = 232$ 年；其平均值为 \qty{220.7}{\text{年}}。

\textbf{2.} 平均错距为 $(7.0 + 6.0 + 5.5 + 7.5)/4 = \qty{6.5}{m}$，且
$T = 6.5/0.030 = \qty{216.7}{\text{年}}$，与平均间隔接近。

\textbf{3.} 最后三次地震的错距为 $6.0 + 5.5 + 7.5 =
\qty{19.0}{m}$，在 1050 至 1712 年间、662 年内累积：$\qty{19.0}{m}
/\qty{662}{\text{年}} = \qty{28.7}{mm.\text{年}^{-1}}$，比
大地测量值 \qty{30}{mm.\text{年}^{-1}} 低 \qty{4}{\percent}。两者在日期与错距的不确定度范围内相符，
该段在其大地震中几乎释放了
其积累的全部加载量。

\textbf{4.} $2025 - 1712 = 313$ 年，亏损量为
$\qty{0.030}{m.\text{年}^{-1}}\times\qty{313}{\text{年}} = \qty{9.39}{m}$.

\textbf{5.} 该亏损量超过\omterm{def:b1:active-margins:active-margin}{海沟}中的每一次错距，其中最大者为
\qty{7.5}{m}，且自 1712 年以来的时间超过了最长间隔 240 年。
在这个意义上，该段已到期；这并不意味着地震
迫在眉睫，只说明断层已积累的滑动量超过了
记录中任何一次地震一次性释放的滑动量。

\textbf{6.} 间隔在 190 至 240 年之间，错距在
\qtyrange{5.5}{7.5}{m} 之间：该循环并非周期性的，下次地震时断层的强度
未知。\omterm{def:b1:active-margins:active-margin}{海沟}资料给出平均值与离散程度，
由此可估计概率，而非具体日期。

\textbf{7.} $A = \qty{250}{km}\times\qty{15}{km} = \qty{3750}{km^2} =
\qty{3.75e9}{m^2}$.

\textbf{8.} $M_0 = \num{3e10}\times\num{3.75e9}\times\num{9.39} =
\qty{1.06e21}{N.m}$.

\textbf{9.} $M_0 = \num{3e10}\times\num{3.75e9}\times\num{7.5} =
\qty{8.44e20}{N.m}$，故 $M_w = \tfrac23(20.926 - 9.1) = 7.88$。

\textbf{10.} 比值为 $9.39/7.5 = 1.25$，相当于震级之差 $\tfrac23\log_{10}1.25 =
0.07$ 个震级。

\textbf{11.} 面积减半使\omterm{def:b1:earthquake-source:moment}{地震矩}减半，震级降低
$\tfrac23\log_{10}2 = 0.20$: $M_w = 7.75$.

\textbf{12.} $\mu A V = \num{3e10}\times\num{3.75e9}\times\num{0.030} =
\qty{3.38e18}{N.m}$ 每年。

\textbf{13.} $b = 1$。$N(5) = 10^{4.3 - 5} = 0.20$（每年），$N(6) = 0.020$
（每年）。

\textbf{14.} $N(7) = 0.0020$（每年）：每 500 年一次；$N(8) = 0.00020$
（每年）：每 \num{5000} 年一次。

\textbf{15.} 该定律给出震级 7.95 或更大的地震每
$1/10^{4.3 - 7.95} = \num{4500}$ 年一次，而\omterm{def:b1:active-margins:active-margin}{海沟}中的地震每 220 年一次，频率高出
二十倍。该定律拟合的是本区所有断层上的大量小地震，
若外推到单个段上的大地震，它
会严重低估它们。某条断层上大地震的复发
应取自该断层自身的记录及其加载情况，而非取自小地震。

\textbf{16.} $\num{5.5e17}/\num{3.38e18} = 0.16$：中等地震
释放了该段所积累\omterm{def:b1:earthquake-source:moment}{地震矩}的 \qty{16}{\percent}。

\textbf{17.} 由\cref{prop:b1:earthquake-source:moment-steps}，
$10^{1.5\times(7.95 - 6)} = 10^{2.92} \approx 840$ 倍。

\textbf{18.} 取 $b = 1$，\omterm{def:b1:earthquake-source:moment}{地震矩}主要由最大的地震决定：
无论中等地震如何，它们释放的只是一小部分，此处为
\qty{16}{\percent}，其余 \qty{84}{\percent} 即滑移亏损，
则等待大地震来释放。

\textbf{19.} 对每一对有 $v(-x) = -v(x)$，故剖面是反对称的。
远离断层处，速度趋于 $\pm V/2 = \pm\qty{15}{mm.\text{年}^{-1}}$。

\textbf{20.} $D = x/\tan(\pi v/V) = 15/\tan(\pi\times7.5/30) = 15/\tan\ang{45} =
\qty{15}{km}$：在\omterm{def:b1:earthquake-source:locking-depth}{闭锁深度}处速度为 $V/4$。

\textbf{21.} 在 \qty{5}{km} 处：$D = 5/\tan(\pi\times3.1/30) = 5/\tan\ang{18.6} =
\qty{14.9}{km}$；在 \qty{40}{km} 处：$D = 40/\tan\ang{69.6} = \qty{14.9}{km}$。
速度精确到 \qty{0.1}{mm.\text{年}^{-1}}，仅这一取整就使
$D$ 约变化 \qty{0.1}{km}；若用实际数据，每个台站的噪声
会更大，\omterm{def:b1:earthquake-source:locking-depth}{闭锁深度}应同时对全部台站进行拟合。

\textbf{22.} \omterm{def:b1:earthquake-source:locking-depth}{闭锁深度} \qty{15}{km}，等于最深地震的
深度。断层的闭锁部分即其脆性部分，即
\omterm{def:b1:stress-strain-rheology:strength-envelope}{孕震层}，摩擦使其在两次地震之间保持锁定，并在地震中破裂；在其下方，
较热的岩石发生蠕动（见\cref{ch:b1:stress-strain-rheology}），断层
稳定地滑动，即模型中以速率 $V$ 滑动、位于 $D$ 以下的部分。

\textbf{23.} $V/(\pi D) = \num{0.030}/(\pi\times\num{15000}) = \num{6.4e-7}$
每年。

\textbf{24.} 取 $M_0 = \qty{1.056e21}{N.m}$，$\log_{10}M_0 = 21.024$，且
$M_w = \tfrac23(21.024 - 9.1) = 7.949$。精确到 0.01，能释放该段
全部亏损量的地震，其\omterm{def:b1:earthquake-source:mw}{矩震级}为
\textbf{7.95}.
\end{solution}
